\author{}
\documentclass[11pt,a4paper]{article}
\usepackage[utf8]{inputenc}
\usepackage[T1]{fontenc}
\usepackage[french]{babel}

% Paquets pour les mathématiques
\usepackage{amsmath, amssymb, array}

% Marges
\usepackage{geometry}
\geometry{hmargin=2cm, vmargin=2cm}

% Paquet pour les boîtes colorées
\usepackage[many]{tcolorbox}

% Définition d'un style personnalisé pour les définitions (bleu)
\newtcolorbox{boiteDef}[1]{
    colback=blue!5!white,      % Couleur de fond (très clair)
    colframe=blue!75!black,    % Couleur du contour et du titre
    title={\textbf{#1}},       % Le titre en gras
    arc=3mm,                   % Arrondi des coins
    boxrule=1pt,               % Épaisseur du trait
    fonttitle=\bfseries        % Police du titre
}

% Définition d'un style pour les propriétés (rose)
\newtcolorbox{boiteProp}[1]{
    colback=pink!5!white,
    colframe=purple!70!pink, 
    title={\textbf{#1}},
    arc=3mm,
    boxrule=1pt,
}

% Définition d'un style pour les rappels (vert)
\newtcolorbox{boiteRappel}[1]{
    colback=green!10!white,
    colframe=green!50!black,  
    title={\textbf{#1}},
    arc=3mm,
    boxrule=1pt,
}

% Définition d'un style pour les remarques (jaune)
\newtcolorbox{boiteRemarque}[1]{
    colback=yellow!10!white,
    colframe=yellow!70!black,  
    title={\textbf{#1}},
    arc=3mm,
    boxrule=1pt,
}

\renewcommand{\thesection}{\Roman{section}}

\begin{document}

\title{\textbf{Chapitre 5 : Limites de fonctions }}
\date{}
\maketitle

\section{Limites en $+\infty$ ou en $-\infty$}

\subsection{Limite infinie en $+\infty$ ou en $-\infty$}

\begin{boiteDef}{Définitions :}
\begin{itemize}
    \item Une fonction $f$ a pour limite $+\infty$ en $+\infty$ si tout intervalle de la forme $]A,+\infty[$, avec $A\in\mathbb{R}$, contient toutes les valeurs $f(x)$ pour $x$ suffisamment grand.
    \newline
    On note alors $\displaystyle \lim_{x \to +\infty} f(x) = +\infty$ ou $f(x) \to +\infty$.
    \item Une fonction $f$ a pour limite $-\infty$ en $+\infty$ si tout intervalle de la forme $]-\infty;A[$, avec $A\in\mathbb{R}$, contient toutes les valeurs $f(x)$ pour $x$ suffisamment grand.
    \newline
    On note alors $\displaystyle \lim_{x \to +\infty} f(x) = -\infty$ ou $f(x) \to -\infty$.
\end{itemize}
\end{boiteDef}

\begin{boiteProp}{Propriété des fonctions de référence :}
    Les fonctions $x\mapsto \sqrt{x}$, $x\mapsto x$, $x\mapsto x^p$ avec $p\in\mathbb{N}^*$ et $x\mapsto e^x$ ont pour limite $+\infty$ en $+\infty$.
    \newline
    En $-\infty$, résultats analogues :
    \begin{itemize}
        \item $\displaystyle \lim_{x\to -\infty} x^n = +\infty$ si $n$ est pair 
        \item $\displaystyle \lim_{x\to -\infty} x^n =-\infty$ si $n$ est impair.
    \end{itemize}
\end{boiteProp}

\begin{boiteRemarque}{Remarque :}
    $x\mapsto \sqrt{x}$ n'est pas définie en $-\infty$ donc cela n'a pas de sens de chercher sa limite en $-\infty$.
\end{boiteRemarque}

\subsection{Limite finie en $+\infty$ ou en $-\infty$}

\begin{boiteDef}{Définition :}
    Une fonction $f$ a pour limite $l\in\mathbb{R}$ en $+\infty$ si tout intervalle ouvert contenant $l$ contient toutes les valeurs $f(x)$ pour $x$ suffisamment grand.
    \newline
    On note alors $\displaystyle \lim_{x \to +\infty} f(x) = l$ ou $f(x) \to l$.
\end{boiteDef}

\begin{boiteProp}{Propriétés des fonctions de référence :}
    Les fonctions $x\mapsto \dfrac{1}{x^n}$ avec $n\in\mathbb{N}^*$ et $x\mapsto \dfrac{1}{\sqrt{x}}$ ont pour limite $0$ en $+\infty$.
    \newline
    Les fonctions $x\mapsto \dfrac{1}{x^n}$ avec $n\in\mathbb{N}^*$ ont pour limite $0$ en $-\infty$.
    \newline
    Les fonctions $x\mapsto e^x$ ont pour limite $0$ en $-\infty$.
\end{boiteProp}

\begin{boiteDef}{Définition :}
    Dans un repère orthogonal, lorsque $\displaystyle \lim_{x\to +\infty} f(x) = l$, alors la courbe représentative de $f$ admet une asymptote horizontale d'équation $y=l$ au voisinage de $+\infty$.
    \newline
    Lorsque $\displaystyle \lim_{x\to -\infty} f(x) = l$, alors la courbe représentative de $f$ admet une asymptote horizontale d'équation $y=l$ au voisinage de $-\infty$.
\end{boiteDef}

\section{Limites en un réel a}

\begin{flushleft}
    $f$ est une fonction définie sur un ensemble $D_{f}$ (un intervalle ou une réunion d’intervalles).
    On peut étudier la limite de $f$ en un réel $a$ si $a \in $ $D_{f}$ ou lorsque $f$ est définie au voisinage de $a$ ($a$ est une borne
    de $D_{f}$).
\end{flushleft}

\subsection{Limite infine en un réel $a$}

\begin{boiteDef}{Définition :}
    Une fonction $f$ a pour limite $+\infty$ en $a$ si tout intervalle de la forme $]A,+\infty[$, avec $A\in\mathbb{R}$, contient toutes les valeurs $f(x)$ pour $x$ suffisamment proche de $a$.
    \newline
    On note alors $\displaystyle \lim_{x \to a} f(x) = +\infty$ ou $f(x) \to +\infty$.
\end{boiteDef}

\begin{boiteRemarque}{Remarques :}
    Graphiquement, on définit de manière analogue une limite $-\infty$ en $a$.
\end{boiteRemarque}

\begin{boiteDef}{Définition :}
    Dans un repère orthogonal, lorsque $\displaystyle \lim_{x\to a} f(x) = \infty$, alors la courbe représentative de $f$ admet une asymptote verticale d'équation $x=a$ au voisinage de $a$.
\end{boiteDef}

\subsection{Limite finie en un réel $a$}

\begin{boiteDef}{Définition :}
    Une fonction $f$ a pour limite $l\in\mathbb{R}$ en $a$ si tout intervalle ouvert contenant $l$ contient toutes les valeurs $f(x)$ pour $x$ suffisamment proche de $a$.
    \newline
    On note alors $\displaystyle \lim_{x \to a} f(x) = l$ ou $f(x) \to l$.
\end{boiteDef}

\begin{boiteRemarque}{Remarque :}
    Si $a \in D_f$ et si $f$ admet une limite $l$ en $a$, alors $f(a) = l$.
\end{boiteRemarque}

\section{Opérations sur les limites}

\subsection{Somme, produit, quotient}

Tous les tableaux vus dans les limites des suites restent valables pour les fonctions.

\subsection{Composition}

\begin{boiteProp}{Propriété :}
    Si $\displaystyle \lim_{x\to a} f(x) = b$ et si $\displaystyle \lim_{X\to b} g(X) = c$, alors $\displaystyle \lim_{x\to a} g(f(x)) = c$.
\end{boiteProp}

\subsection{Méthodes et exemples types}

Limites à l'infini : on factorise par ce qui semble être le plus fort.
Pour la forme $\dfrac{0}{0}$, pour lever l'indétermination, on factorise (avec delta ou autre).

\section{Limites et comparaison}

\subsection{Limite infine}

\begin{boiteProp}{Théorème de comparaison :}
    Pour $x$ suffisamment grand, si $f(x) \leq g(x)$ et si $\displaystyle \lim_{x\to +\infty} f(x) = +\infty$, alors $\displaystyle \lim_{x\to +\infty} g(x) = +\infty$.
    \newline
    De même, pour $x$ suffisamment grand, si $f(x) \leq g(x)$ et si $\displaystyle \lim_{x\to +\infty} g(x) = -\infty$, alors $\displaystyle \lim_{x\to +\infty} f(x) = -\infty$.
\end{boiteProp}

\subsection{Limite finie}

\begin{boiteProp}{Théorème des gendarmes :}
    Si $f$, $g$ et $h$ sont trois fonctions et $l$ un nombre réal tel que :
    \begin{itemize}
        \item Pour $x$ assez grand, $f(x) \leq g(x) \leq h(x)$
        \item $\displaystyle \lim_{x\to a} f(x) = l$,
        \item $\displaystyle \lim_{x\to a} h(x) = l$,
    \end{itemize}
    alors $\displaystyle \lim_{x\to a} g(x) = l$.
\end{boiteProp}

\subsection{Croissances comparées}

\begin{boiteProp}{Théorème des croissances comparées :}
    $\displaystyle \lim_{x\to +\infty} \dfrac{e^x}{x^n} = +\infty$ pour tout $n\in\mathbb{N}$.
    \newline
    $\displaystyle \lim_{x\to -\infty} x^n e^x = 0$ pour tout $n\in\mathbb{N}$.
    \newline
    En particulier :
    \newline
    $\displaystyle \lim_{x\to +\infty} \dfrac{e^x}{x} = +\infty$.
    \newline
    $\displaystyle \lim_{x\to -\infty} x e^x = 0$.
\end{boiteProp}


\end{document}